\section{The residual full-owner primitive transversal, directly}
\label{sec:primitive-transversal}
\label{sec:consolidated-owner}

We work after removal of nonidentity finite-class exact kernels,
with a minimal infinite exact $E$ satisfying the \emph{residual}
condition $R_E=F=\Sigma(\Id)$.  The finite regular thin
translation group is $G\cong C_2^d$, and $H=E\cap F$ is its
subgroup preserving an actual $E$-owner.  Let
$B\subseteq M/E$ be a nonempty finite $G$-invariant induced
owner substructure.  Put $N_B=\pi_E^{-1}(B)$, with all these
actual owners named, and $\Gamma=\Aut(B)$ in its actual
quotient orbital language.  No root frame or selected section
is assumed.

\begin{theorem}[Residual real closure over named owners]
\label{thm:relative-real-closure}
For every finite $A\subseteq N_B$ one has
\begin{equation}\label{eq:full-owner-acl}
 \operatorname{acl}_{N_B}(A)=
 \operatorname{acl}^{\rm eq}_{M}(A\cup B)\cap N_B=F[A].
\end{equation}
Moreover $F[A]$ is definable in $N_B$ from $A$ via the
literal named thin neighbour translations.
\end{theorem}
\begin{proof}
Since $R_E=F$, $E$ is the \emph{unique} minimal
infinite invariant equivalence: for a distinct minimal
$E'$, the exact correspondence calculation gives
$E'\subseteq R_E=F$, impossible for infinite $E'$.
There is now a SHORT second source route
to the crucial real closure equality:
\cref{cor:rh-acl-bounded-character}
proves, from the original canonical real
hull and full invariant lattice, that
when the number of minimal infinite
directions is ONE, real algebraic
closure has character one and
$\operatorname{acl}_M(\{a\})=F[a]$.
Consequently
$\operatorname{acl}_M(A)=F[A]$
for every finite real $A$, including
the empty-base convention in this
infinite transitive source.
The earlier direct real-hull argument
\cref{lem:consolidated-unary-closure}
is an independent older route, not a
further assumption needed here.  By the harmless-owner
criterion, Lemma~\ref{lem:acf-harmless-owners}, naming
any finite set of $E$-classes leaves this real closure
unchanged.  Every $t_g$ is a named 0-definable functional
thin edge, so $F[A]$ is definable over $A$ and stays in
$N_B$ by the $G$-invariance of $B$.

To compute real algebraic closure in the owner-named
structure, no joint lifting or homogeneity of $N_B$
is needed. Let
$G_{A,B}=\Aut(M/A,B)$, with $B$ treated as the
ACTUAL finite tuple/set of named original exact
$E$-quotient imaginary classes. Every element
of $G_{A,B}$ preserves the union $N_B$, all
literal original fine colours, every named
owner predicate and every $t_g$. Hence restriction
gives a subgroup of $\Aut(N_B/A)$.
For $x\in N_B\setminus F[A]$,
the equality $\operatorname{acl}_M(A)=F[A]$
and the harmless-owner lemma show that
$x\notin\operatorname{acl}^{\rm eq}_M(A\cup B)$.
Its $G_{A,B}$-orbit is therefore INFINITE,
lies entirely in $N_B$, and is contained
in its $\Aut(N_B/A)$-orbit. Thus
$x\notin\operatorname{acl}_{N_B}(A)$.
Conversely each point of $F[A]$ is a value
of one of the finite named total thin
translations $t_g$ on some $a\in A$,
hence is definable over $A$ in $N_B$.
This proves both displayed algebraic-closure
equalities, including $A=\varnothing$.

SEPARATELY, the owner-named union $N_B$
is homogeneous, but this requires one
MIXED-ORBIT check before joint lifting.
For an original fine atom $a$ and the
same-sort invariant exact owner kernel
$E$, original symmetry makes all
invariant represented relations commute,
so
\[
              \Id;a;E=a;E=E;a;E .
\]
By \cref{lem:nested-mixed}, the mixed
ACTUAL $\Aut(M)$-orbital of one REAL
point $x$ and an $E$-owner class $C$
therefore depends ONLY on the quotient
pair orbital of $([x]_E,C)$, not on
which representative $x$ was chosen
inside its owner class. Thus every
finite owner-preserving fine-colour
partial isomorphism of $N_B$,
extended by the identity on ALL
the finite named $E$-classes $B$,
really is a finite partial
isomorphism of the joint
original REAL/EXACT E-sort language.
Joint lifting
\cref{thm:typed-joint-exact-lifting}
now supplies one ambient
automorphism fixing $B$.
The resulting automorphism preserves
$N_B$, and its restriction extends the
original finite owner-named map.
This extra homogeneity statement is used
in subsequent owner-frame constructions,
but is NOT a prerequisite for the preceding
stabiliser-orbit algebraicity argument.
No group-theoretic section has been assumed.
\end{proof}

\begin{lemma}[Finite avoidance over a full-packet closed owner base]
\label{lem:relative-finite-avoidance}
Let $C\subseteq A\subseteq N_B$ be finite full-$F$-closed
induced structures, and let $Z\subseteq N_B\setminus C$ be finite.
There exists an owner-preserving $g\in\Aut(N_B/C)$ with
\[
                   g(A\setminus C)\cap Z=\varnothing.
\]
\end{lemma}
\begin{proof}
Take the point action of $K=\Aut(N_B/C)$ on $N_B$.
Every $a\in A\setminus C$ has an infinite $K$-orbit,
since $C=F[C]=\operatorname{acl}_{N_B}(C)$ by
\cref{thm:relative-real-closure}.  Apply
\cref{lem:finite-action-avoidance} to the finite moving
set $A\setminus C$ and the finite forbidden set $Z$.
The resulting $g\in K$ fixes $C$ and moves the ENTIRE
finite off-base set simultaneously as required.
\end{proof}

\begin{corollary}[Strong amalgamation of full packets]
\label{cor:relative-closed-sap}
The finite induced owner-typed structures in $N_B$ that
are closed under $F$ strongly amalgamate over their
full-$F$-packet closed common bases, with output again
full-packet closed.
\end{corollary}
\begin{proof}
For a finite full-$F$-closed base $A$, every point of
an extension outside $A$ has infinite orbit under the
automorphisms fixing $A$ and all named owners, by the
preceding theorem.  Apply the finite-coset avoidance
lemma to move all finitely many off-base points of one
factor away from those of the other, simultaneously.
Take the induced union and then its finite $F$-closure;
this closure does not identify any distinct old points or
leave $N_B$, since $F$ is a named finite packet relation.
The two old factor images intersect exactly in $A$.
\end{proof}

\begin{theorem}[The all-owner primitive transversal]
\label{thm:joint-transversal}
\label{lem:full-owner-transversal}
Let $\mathcal U_B$ comprise finite owner-typed diagrams
realised in $N_B$ on vertices from pairwise distinct
full $F$-packets.  Its pair palettes are
\[
          P_{st}=L_{st}\setminus F\qquad(s,t\in B),
\]
with only diversity colours between distinct formal points,
even for repeated owner names.  Then $\mathcal U_B$ is
triangle-defined, has JEP and SAP, and satisfies semantic
homomorphic CCO with injective packet-separated completions.
Every individual owner sort is primitive.  The Fra\"isse
limit $U_B$ recovers $N_B$ under full $G$-packet inflation:
\begin{equation}\label{eq:full-owner-inflation}
 \operatorname{own}(x,g)=g\operatorname{own}(x),\qquad
 \col((x,g),(y,h))=(g+h)\col(x,y)\quad(x\ne y).
\end{equation}
The group $\Gamma=\Aut(B)$ permutes only the owner tags
and fixes the literal raw fine-colour names.
\end{theorem}
\begin{proof}
Any partial transversal diagram whose specified cycles
have typed realisations has a raw homomorphic completion
by semantic CCO of $M$.  Each typed edge is compatible
with its prescribed actual quotient-owner pair orbital.
Anchored relative completion,
Theorem~\ref{thm:anchored-relative}, restores all prescribed
$E$-owners.  The resulting homomorphism may identify
unjoined formal vertices or place them in one $F$-packet.
Resolve collisions one point at a time: if $x$ has a
specified neighbour set $C$ (images of those neighbours),
all its displayed edge colours lie outside $F$, so the
image of $x$ is not in $F[C]$.  By~\eqref{eq:full-owner-acl}
its owner-preserving orbit over $C$ is infinite and hence
meets infinitely many $F$-packets.  Move only the image
of $x$ into a packet disjoint from the finitely many
currently used packets, preserving all specified incident
colours and the actual owner label.  A complete protected
base is fixed by moving only the colliding off-base point.
After finitely many steps the map is packetwise injective.

Every finite complete triangle-legal typed diagram admits
a raw embedding by triangle fullness of $M$, and the
anchored theorem then puts it over the prescribed owners.
The same packet separation makes it a member of
$\mathcal U_B$, proving triangle-definedness.  To strongly
amalgamate transversal diagrams, inflate both to their
entire full $G$-packets, apply the preceding closed-SAP
corollary and select one representative of every packet.
The same works over the empty base, proving JEP.
The usual Fra\"isse theorem gives the homogeneous $U_B$.

We check primitivity of each named owner sort.  Suppose
$D\subsetneq P_{ss}=E\setminus H$ is the nonempty
diversity support of a nontrivial invariant equivalence
on owner sort $s$.  Then $D*D\subseteq D$ in the strong
SAP algebra of distinct transversal vertices.
For $g\in H$, if $D\cap gD\ne\varnothing$, choose
$a,ga\in D$.  Given any $b\in D$, choose
$c\in a*b$; SAP makes this nonempty, and $c\in D$.
Translate the first endpoint of the triangle to obtain
$(ga,b,gc)$; hence $gc\in D$.  Translate instead its
third endpoint to obtain $(a,gb,gc)$; hence $gb\in D$.
Therefore $gD\subseteq D$, and involutivity gives equality.
Thus $H_D=\{g\in H:gD=D\}$ is a subgroup.
Every product of two $D$-atoms outside $H$ lies in $D$;
if it produces a thin $t_g$ then $D\cap gD\ne\varnothing$,
so $g\in H_D$.  The raw relation
\[
               \Id\cup\{t_g:g\in H_D\}\cup D
\]
is consequently an invariant equivalence contained in $E$.
Here the invariance concerns the ORIGINAL
\(\Aut(M)\) action, not only the named-owner
sort action: \(D\) is a union of LITERAL old
fine atomic orbitals and each \(t_g\) is an old
globally named thin orbital. Thus this RAW
relation is a uniform old-colour union on all
of \(M\), with NO chosen owner section.
It is also globally an EQUIVALENCE:
any pair of consecutive related edges
lies inside one actual \(E\)-class.
Transitivity of \(M\) moves that class
to the chosen owner \(s\), where the
preceding thin/nonthin product check
established closure under composition.
All these checks are in the original
old fine-colour language; they do not
postulate a new named local relation.

It has INFINITE classes already by
the SAP of the OWN transversal,
without another appeal to the
canonical real-hull theorem.
Indeed choose a genuinely realised
\(D\)-edge \(xy\) on the single named
owner sort \(s\). For every
integer \(n\ge1\), strongly
amalgamate \(n\) copies of this
COMPLETE two-point transversal
diagram over the SAME singleton
\(x\). SAP preserves every
\(D\)-edge from \(x\) and gives
\(n\) pairwise DISTINCT
\(D\)-neighbours \(y_1,\ldots,y_n\),
lying in distinct full \(F\)-packets.
The resulting finite diagram
is in the already established
age \(\mathcal U_B\); by the
DEFINITION of that age it has
an ACTUAL realisation in \(N_B\),
and therefore in original \(M\).
All these \(D\)-neighbours lie
in ONE class of the displayed
original global equivalence,
because each original \(D\)-edge
belongs to its named diversity
support. Thus for every \(n\)
some original global class has
at least \(n+1\) points.
Original transitivity of \(M\)
makes all classes of an
\(\Aut(M)\)-invariant equivalence
equipotent, so EVERY such class
is infinite. This argument
does not require a FIXED
owner-preserving real embedding
of the infinite transversal,
nor identify arbitrary
unrecorded pair colours.

Since \(E\) is the unique
minimal infinite invariant
equivalence in this residual
source and the original
invariant-equivalence lattice
is FINITE, every invariant
equivalence with infinite
classes contains \(E\).
The displayed equivalence
is contained in \(E\);
hence it equals \(E\).
On the original owner sort,
its nonthin diversity part
is \(D\), while that of \(E\)
is \(E\setminus F\).
Thus equality forces
\(D=E\setminus H\),
contradicting properness.

Finally, complete finite $G$-packet diagrams in $N_B$
and full inflations of finite $\mathcal U_B$ diagrams
have the same induced fine colours, by the thin transport
law.  For any finite full-packet embedding and extension,
choose origins on its packets, apply the Fra\"isse
extension property to the resulting transversal map,
and recover every packet point by its unique $G$-translation
from that origin.  This preserves all fine matrices and
owner labels, with no equivariant choice of origins.
Thus the inflated generic is homogeneous with the same age
as the homogeneous $N_B$, proving the claimed isomorphism.
The pure $\Gamma$ language action comes from ambient exact
owner homogeneity and preserves every literal raw colour.
\end{proof}

\begin{remark}
The theorem supplies exactly the typed SAP, primitivity,
closed packet completion and owner $\Gamma$ actions used by
the selector/HKN construction.  The more general conditional
bijection groupoid and root frame geometry remain available in
the original R45 sources but are not premises of the
consolidated residual construction.
\end{remark}
